% =============================================================================
%  BSc Thesis — University of Debrecen
%  Enhancing Workforce Scheduling and Operational Efficiency
%  Using AI-Powered Predictive Analytics
%  Author: Maua Elaine Imani
%  Supervisor: Dr. Bogacsovics Gergő, Assistant Professor
% =============================================================================

\documentclass[12pt,a4paper,oneside]{article}

\usepackage{mathptmx}
\usepackage{xspace}
\usepackage[T1]{fontenc}
\usepackage[utf8]{inputenc}
\usepackage{textcomp}
\usepackage{setspace}
\usepackage{caption}
\captionsetup{
  font={small,it},
  labelfont={small,bf},
  labelsep=period,
  justification=centering
}
\onehalfspacing
\usepackage[left=3cm, right=2cm, top=3cm, bottom=3cm]{geometry}
\usepackage{xcolor}
\definecolor{thesisblue}{RGB}{0, 83, 156}
\definecolor{linkblue}{RGB}{0, 102, 204}
\definecolor{citegreen}{RGB}{0, 128, 0}
\definecolor{codegray}{RGB}{245, 245, 245}
\definecolor{codetext}{RGB}{40, 40, 40}
\usepackage{graphicx}
\graphicspath{{figures/}}
\usepackage{subcaption}
\usepackage{float}
\usepackage{tikz}
\usetikzlibrary{shapes.rounded,arrows.meta,positioning}
\usepackage{booktabs}
\usepackage{amsmath,amssymb,amsfonts}
\usepackage{amsthm}
\usepackage{listings}
\lstset{
    basicstyle=\small\ttfamily\color{codetext},
    breaklines=true,
    frame=single,
    framerule=0.5pt,
    rulecolor=\color{gray!50},
    backgroundcolor=\color{codegray},
    captionpos=b,
    numbers=left,
    numberstyle=\tiny\color{gray},
    numbersep=8pt,
    tabsize=4,
    showstringspaces=false
}
\usepackage{enumitem}
\setlist{nosep, leftmargin=*}
\setlist[itemize]{label=\textbullet}
\setlist[enumerate]{label=\arabic*.}
\usepackage{titlesec}
\titleformat{\section}{\fontsize{14}{17}\bfseries}{\thesection}{1em}{}
\titleformat{\subsection}{\fontsize{13}{16}\bfseries}{\thesubsection}{1em}{}
\titleformat{\subsubsection}{\normalsize\bfseries\itshape}{\thesubsubsection}{1em}{}
\titlespacing*{\section}{0pt}{2.5ex plus 1ex minus .2ex}{1.5ex plus .2ex}
\titlespacing*{\subsection}{0pt}{2ex plus 1ex minus .2ex}{1ex plus .2ex}
\titlespacing*{\subsubsection}{0pt}{1.5ex plus 1ex minus .2ex}{0.5ex plus .2ex}
\usepackage[authoryear,round]{natbib}
\bibliographystyle{plainnat}
\usepackage[
    colorlinks=true,
    linkcolor=thesisblue,
    citecolor=citegreen,
    urlcolor=linkblue,
    bookmarks=true,
    bookmarksnumbered=true,
    pdfstartview=FitH,
    pdftitle={Enhancing Workforce Scheduling and Operational Efficiency Using AI-Powered Predictive Analytics},
    pdfauthor={Maua Elaine Imani}
]{hyperref}
\widowpenalty=10000
\clubpenalty=10000
\hyphenpenalty=10000
\exhyphenpenalty=10000
\tolerance=9999
\emergencystretch=3em
\setlength{\parindent}{1.25cm}
\setlength{\parskip}{0pt}
\usepackage{fancyhdr}
\pagestyle{fancy}
\fancyhf{}
\fancyhead[L]{}
\fancyhead[R]{}
\fancyfoot[C]{\small\thepage}
\renewcommand{\headrulewidth}{0pt}
\renewcommand{\footrulewidth}{0pt}
\fancypagestyle{plain}{
    \fancyhf{}
    \fancyfoot[C]{\small\thepage}
    \renewcommand{\headrulewidth}{0pt}
}
\newcommand{\eg}{e.g.,\xspace}
\newcommand{\ie}{i.e.,\xspace}
\newcommand{\etal}{et al.\xspace}

\begin{document}

%  --  EXTERNAL COVER  --
\begin{titlepage}
\thispagestyle{empty}
\null\vfill
\begin{center}
{\fontsize{18}{24}\selectfont\bfseries THESIS}
\end{center}
\vfill
\hfill
\begin{minipage}[t]{0.4\textwidth}
\raggedleft
\textbf{Maua Elaine Imani}
\end{minipage}
\vfill\vfill
\begin{center}{\large Debrecen, 2026}\end{center}
\vspace{1cm}
\end{titlepage}

\newpage

%  --  INTERNAL TITLE PAGE  --
\begin{titlepage}
\thispagestyle{empty}
\begin{center}
{\Large\scshape University of Debrecen}\\[0.3cm]
{\large\scshape Faculty of Informatics}\\[0.2cm]
{\scshape Department of Data Science and Visualization}\\[1.8cm]
{\fontsize{16}{20}\selectfont\bfseries
Enhancing Workforce Scheduling and\\[0.3cm]
Operational Efficiency Using\\[0.3cm]
AI-Powered Predictive Analytics}
\end{center}
\vspace{6.8cm}
\begin{minipage}[t]{0.48\textwidth}
\begin{flushleft}
\textbf{Supervisor:}\\[0.2cm]
Dr.~Gerg\H{o} Bogacsovics\\
\textit{Assistant Professor}
\end{flushleft}
\end{minipage}
\hfill
\begin{minipage}[t]{0.48\textwidth}
\begin{flushright}
\textbf{Author:}\\[0.2cm]
Maua Elaine Imani\\
\textit{Computer Science BSc}
\end{flushright}
\end{minipage}
\vspace{1.8cm}

\begin{center}
%\textit{``I consent to the submission of this work for review.''}
\end{center}

\vfill

\begin{center}
{\large Debrecen, 2026}
\end{center}

\end{titlepage}

%  --  TABLE OF CONTENTS  --
\newpage
\thispagestyle{plain}
\tableofcontents
\newpage

% =============================================================================
% SECTION 1 — INTRODUCTION
% =============================================================================
\section{Introduction}
\label{sec:introduction}

%\subsection*{Background and Motivation}

To the naked eye, creating a weekly schedule is a simple task, almost like a simple for loop. For each employee, for each shift, if a shift slot is empty, assign employee to shift. However, scheduling the right number of employees for a week is a brittle task that involves a lot more variables --  demand changes day to day, employees are bound by certain legal working hours, labor laws enforce additional constraints, and shifts need to be distributed fairly across a team. There are also sudden happenings like cancellations and large last-minute bookings that the same schedule needs to adapt to. The needs of a restaurant on Monday may look very different from Tuesday. In fact, the needs of a Monday morning may look very different from the same Monday's afternoon. Despite all this, most hospitality businesses still schedule by hand -- a fixed rota adjusted from memory and experience \citep{thompson2004}.

Manual scheduling is incapable of responding to the full range of demand signals that drive customer volumes. An experienced manager could possibly make good calls regarding an approaching holiday or an in-house event.  They could even reasonably infer that since occupancy of the larger hotel is high, then the restaurant outlet would also be busy. It becomes a completely different matter when they have to simultaneously integrate hotel occupancy rates, daily temperature forecasts, day-of-week patterns, and seasonal trends into a coherent staffing decision. Doing so across five distinct roles for two or three shift periods in a given day, seven days a week, for all the employees in the pool exceeds what any single person, no matter the experience, can reliably handle. In practice, this means that a manager schedules the same headcount week after week, regardless of whether the hotel is nearly full or sitting at 25\% occupancy \citep{choi2009}.

This demand-blind scheduling is costly both for the team on the floor and the business. Overstaffing on low-demand days means that there are people on the floor being paid yet they have no reasonable impact on the business. And for the staff, it often means being sent home early with fewer hours than expected. Understaffing on busy days cuts the other way: the team is overwhelmed, service quality drops, and covers get lost when guests cannot be seated or served adequately \citep{choi2009}. Poor service quality compounds the problem further. A guest who had a bad experience is unlikely to return, and if they leave a review, one who reads those reviews may never visit at all \citep{sparks2011}. These outcomes are completely avoidable. The demand signals that drive these numbers are observable weeks or days in advance. For instance, hotel occupancy forecasts are available through the property management system (PMS), temperature forecasts are freely available, and holidays are deterministic. The challenge is then not the information but a mechanism that translates that information into staffing decisions.

This thesis presents \textit{SlotMeIn}, a scheduling engine designed for hotel food and beverage operations. The engine addresses the scheduling challenge, first using a machine learning component, then with a constraint-satisfier scheduler. The machine learning component is trained on real operational data and learns patterns between the demand signals to predict a tier, either Low, Medium, or High, for each shift period on each day. The constraint satisfier scheduler then uses these predicted demand tiers to determine the headcount and generate a compliant assignment that reflects the needs of a given day.

The thesis uses exploratory analysis of real covers data from a Budapest city hotel restaurant to demonstrate that the afternoon cover volumes are significantly predicted by hotel occupancy rates, mean daily temperature, day of the week, season of the year and the presence of Hungarian public holidays. An ordinary least squares regression fitted to the 294 days of observed cover counts  yields an adjusted $R^2$ of 0.43. This confirms that nearly half of the variance in demand is explained by observable demand signals before the restaurant doors even open.

In the study, four candidate classifiers, that is, a rule-based occupancy baseline, Logistic Regression, Random Forest, and XGBoost, are trained and evaluated on the last 20\% of every season's data. Logistic Regression performs the best at 62.5\% and is selected among the four. The dominance of Logistic Regression over the ensemble methods suggests that the demand-tier relationship is largely linear, which aligns with the OLS findings.

Initially, the scheduling algorithm that was selected was greedy, before it was replaced by a constraint programming-satisfiability (CP-SAT) solver \citep{ortools2024}. The CP-SAT solver has abilities such as backtracking and optimizing for several objectives using a single weighted function that a greedy algorithm simply does not have. The move from greedy to CP-SAT concretely shows why greedy falls short when the problem space gets complex as in scheduling.

SlotMeIn shows that the information needed to generate dynamic schedules already exists, it just needs a system that knows how to churn it.

\newpage

% =============================================================================
% SECTION 2 — LITERATURE REVIEW
% =============================================================================
\section{Literature Review}
\label{sec:literature}

This literature review section examines the theoretical foundations and existing approaches relevant to the development of an intelligent, demand-aware scheduling engine for Food \& Beverage operations within the hospitality industry.

\subsection{Workforce Scheduling in Hospitality}

Workforce scheduling in hospitality involves a multi-step process that begins long before a single assignment is made. \citet{thompson2004} breaks this down into four sequential phases:

  \begin{enumerate}
  \item Forecasting demand
  \item Translating forecasts into staffing requirements
  \item Scheduling employees to meet those requirements
  \item Controlling the live plan on the day
  \end{enumerate}

Workforce scheduling has two types of constraints -- hard and soft. \citet{thompson2004} gives the differences between hard and soft constraints. Hard constraints are invariants that cannot be violated by the final schedule at all. These include maximum daily working hours, mandatory rest periods and role eligibility requirements. Soft constraints are desirable but not mandatory to be satisfied. For instance, if an employee prefers a Monday morning shift over a Thursday evening shift, that is considered a soft constraint. This distinction shapes the design of the CP-SAT engine described in Section~\ref{sec:scheduler}.

\citet{ernst2004} survey over 700 papers published between 1950 and 2003, documenting  fields such as healthcare, transportation, retail, and hospitality in which scheduling has been studied. One key takeaway is that even moderately sized scheduling problems are NP-hard \citep{ernst2004}. NP-hard means that as a problem size grows, the computational time required to find an optimal solution grows exponentially. More employees, more shift types, and longer planning horizons all compound this effect. This is why approaches like CP-SAT are used: finding the best solution fast matters more than proving that it is the best one.


\citet{bergh2013} offer a more recent and targeted survey of the shift scheduling literature, distinguishing between cyclic and acyclic scheduling patterns. In cyclic scheduling, the same patterns repeat themselves and there is little variance. For instance, an 8.00 a.m. class is set once at the beginning of a semester and does not depend on external factors such as weather to move the class to a different time slot. Acyclic scheduling, on the other hand, varies by demand from period to period and the schedule must adapt accordingly. The hospitality context studied in this thesis is inherently acyclic. Demand is shaped by seasonal patterns, local events, weather, and public holidays that produce different requirements every week. A cyclic schedule cannot meaningfully capture these variations and is precisely the type of scheduling this thesis aims to replace.


\subsection{Demand Drivers for F\&B Covers}

\citet{thompson2008} tested 27 forecasting methods across 41 outlet-daypart combinations at four hotels, to determine whether hotel room occupancy data improves forecasting accuracy. In 34 out of 41 of those scenarios, the models that incorporated occupancy data outperformed those without \citep{thompson2008}. In all these cases, the accuracy improvement exceeded 11\%. The model that won most frequently was linear regression, winning in 12 out of the 41 scenarios. This directly justifies why occupancy is the primary feature used in this thesis.

Weather conditions, seasonal patterns and special occasions have also been identified as demand drivers in the broader forecasting context \citep{pmc2025, heliyon2023}. The feature set used in this study was built on this evidence and validated against real data through OLS regression.

The weather signal is particularly important in the context of Budapest. Budapest experiences continental climate with seasons whose temperatures vary greatly. The mean daily temperatures in the Open-Meteo records used in this study \citep{openmeteo2023} range from approximately  $-2$°C in January to $+25$°C in July. This range is enough to detect the temperature effects independently of occupancy.  Since temperature and occupancy are partially correlated, both were included in the OLS regression so their individual contributions could be separated.


\subsection{Machine Learning for Staffing Prediction}

Classical time-series methods like ARIMA treat demand as a univariate sequence, with no consideration for the signals \citep{hyndman2018}.  Tree-based ensemble methods like Random Forests \citep{breiman2001} and XGBoost \citep{chen2016} were included as candidate models because of their ability to handle non-linear feature interactions. Logistic Regression was included as a linear baseline.

\subsection{Synthetic Data Generation for Machine Learning}

Data scarcity is a fundamental problem when it comes to training ML models for staff scheduling. The data generated for a given restaurant accumulates slowly. This thesis had access to 294 days of real cover data which is sufficient for regression but not so for robust classifier training \citep{lu2023}. This motivated the move to using synthetic data for model training. \citet{lu2023} identify statistical difference evaluation as an appropriate validation strategy for synthetic tabular data. Kolmogorov--Smirnov test is the standard implementation.

\subsection{Constraint Satisfaction for Employee Scheduling}

The Nurse Scheduling Problem (NSP) is the canonical analogue to restaurant shift scheduling in the operations research literature \citep{bensaid2024}. Nurses become food and beverage staff, wards become shift periods, and the set of hard constraints changes from healthcare to hospitality.


Greedy constructive heuristics are widely used as scheduling baselines since they are fast and simple to implement \citep{bergh2013}. However, they are notoriously myopic. Once a locally optimal decision is made, the greedy algorithm locks it in which closes the door to a more globally correct rendition.

CP-SAT, the constraint programming solver in Google's OR-Tools library \citep{ortools2024}, extends classical constraint programming with efficient SAT-based reasoning. CP-SAT searches the problem space systematically and backtracks when a partial assignment violates hard constraints. This ability is essential when hard constraints interact in ways that make apparently valid assignments infeasible.


\subsection{EU and Hungarian Labor Law in Scheduling}

The hard constraints that are used in SlotMeIn are derived directly from applicable labor law. The EU Working Time Directive \citep{wtd2003} and the Hungarian Labour Code \citep{hungarian2012} establish three non-negotiable labor requirements:

  \begin{itemize}
  \item A minimum of 11 consecutive hours of rest between working days
  \item No more than 48 hours average weekly working time over a reference period
  \item No more than six consecutive working days without a rest day
  \end{itemize}

These are encoded as hard constraints  in the CP-SAT formulation.

\newpage
% =============================================================================
% SECTION 3 — DATA AND METHODOLOGY
% =============================================================================
\section{Data, Simulation, and Methodology}
\label{sec:data}

This chapter presents the data sources, simulation methodology, background for the need of simulation, and the research methods used to develop the proposed demand-aware workforce scheduling system.

\subsection{Overview of Data Sources}


The empirical work in this thesis is founded on four data sources, all coming from or directly related to a Budapest city hotel. Table~\ref{tab:datasources} gives a summary.


\begin{table}[H]
\centering
\caption{Primary data sources used in this thesis.}
\label{tab:datasources}
\small
\renewcommand{\arraystretch}{1.6}
\begin{tabular}{@{}p{3.8cm}p{3.2cm}@{\hspace{0.4cm}}p{2.5cm}@{}}
\toprule
\textbf{Dataset} & \textbf{Period} & \textbf{Records} \\
\midrule
Afternoon covers & \mbox{Jan 2025 -- Apr 2026} & 294 days \\
Hotel bookings & \mbox{Jul 2023 -- Apr 2026} & 119,390 \\
\mbox{Open-Meteo weather} & \mbox{Jul 2023 -- Apr 2026} & 1,008 days \\
Hungarian holidays & 2023--2026 & Deterministic \\
\bottomrule
\end{tabular}
\renewcommand{\arraystretch}{1}
\end{table}

\paragraph{Afternoon covers:} These are the actual headcounts of the guests who came to the restaurant in the afternoon, collected over 15 months.  These were normalized by dividing by the 95th percentile of the distribution, used as a full-house proxy. The resulting rates run from 0.0 to 1.0, with a mean of 0.65 and standard deviation of 0.22.


\paragraph{Hotel bookings:} These are daily occupancy rates computed from 119,390 booking records pulled from the hotel's property management system. Occupancy is expressed as a percentage of total occupancy. It is available to management before the day starts through standard PMS reporting.


\paragraph{Weather data:} These are daily temperature records for Budapest, obtained from the Open-Meteo historical weather API \citep{openmeteo2023}. Temperature is a rough signal for guest outdoor dining behavior and tourism activity. Warm summer days might nudge people to step out for a bite or drink, while extreme cold encourages staying indoors.


\paragraph{Hungarian public holidays:} All eleven official Hungarian public holidays were included as binary flags, along with four dates that are not statutory but drive noticeable demand shifts. These days are: Christmas Eve, New Year's Eve, Valentine's Day, and Easter Sunday. Since Easter Sunday is not a specific calendar date, it was computed using the \texttt{dateutil.easter} algorithm.

These four sources were joined on calendar date to produce a unified feature table of 294 complete observations, each containing the actual afternoon cover count, hotel occupancy percentage, daily temperature readings, day-of-week and weekend indicators, season classification, and the special-day binary flag.


\begin{table}[H]
\centering
\small
\renewcommand{\arraystretch}{1.3}
\begin{tabular}{@{}lrrllrrrr@{}}
\toprule
\textbf{Date} & \textbf{Covers} & \textbf{Occ.\%} & \textbf{Season} & \textbf{Day} & \textbf{WE} & \textbf{Special} & \textbf{Temp mean} & \textbf{Temp max} \\
\midrule
2025-01-10 & 124 & 19.21 & Winter & Friday    & 0 & 0 &  3.8 &  9.5 \\
2025-01-22 & 117 & 29.80 & Winter & Wednesday & 0 & 0 & -0.7 &  0.6 \\
2025-03-03 & 179 & 71.52 & Spring & Monday    & 0 & 0 &  3.3 &  9.1 \\
2025-03-27 & 308 & 50.99 & Spring & Thursday  & 0 & 0 & 12.3 & 17.6 \\
2025-04-27 & 149 & 31.79 & Spring & Sunday    & 1 & 0 & 12.8 & 18.8 \\
\bottomrule
\multicolumn{9}{@{}l}{\small WE = is\_weekend flag (1 = weekend); Temp in °C.}
\end{tabular}
\renewcommand{\arraystretch}{1}
\caption{First five rows of the unified feature table ($n = 294$).}
\label{tab:featurehead}
\end{table}

\subsection{Exploratory Analysis and Data Partitioning}

\subsubsection{Exploratory Data Analysis}
\label{sec:eda}

The EDA was run before any modeling to check whether the features actually relate to cover counts in a visible way. If nothing showed up in the plots, training a classifier on these features would not have made sense.

 Figure~\ref{fig:demand_signals} shows six panels across the 294 real observations. Occupancy and temperature both show positive associations with afternoon covers. As occupancy rises, covers tend to follow. The relationship looks roughly linear across the observed range, with no obvious plateau or cutoff point. Temperature tells a similar story but with more noise, which tracks with the OLS results where the temperature coefficient carries a larger standard error. Some of that noise is expected. Warm days and high occupancy tend to happen at the same time in summer, so the raw scatter mixes the two effects together.

\begin{figure}[H]
    \centering
    \includegraphics[width=1.1\textwidth]{figures/demand_vs_covers.png}
    \caption{Demand signals vs actual afternoon covers.}
    \label{fig:demand_signals}
\end{figure}


Summer (June - August)  consistently shows the highest cover counts. Winter is generally quiet but December has a pattern of its own. Late November ramping up to mid-December follows the winter pattern. However, the second half of the month, leading up to Christmas and New Year's Eve picks back up sharply. This split within December is what motivated the festive ramp in the simulation design (Section~\ref{sec:simulation}).

Hungarian public holidays and operationally significant dates are associated with higher cover counts than normal days. The median is roughly 60 -- 70 covers higher, and the spread is wider because not all special days are equally busy. Weekends show a similar pattern, but at a smaller scale. Saturdays and Sundays run about 20-30 more covers than weekdays.

These patterns confirm that demand is not random. There are observable signals that can be measured in advance and incorporated into a forecasting model to help create demand-aware staffing decisions. 


\subsubsection{Data Partitioning}
\label{sec:split}

The dataset was split using a season-stratified approach. The last 20\% of every season's days were used for testing, leaving the other 80\% for  training. Splitting by season means both sets saw patterns across all four seasons. A naive chronological split would concentrate test data squarely in the winter and early spring, giving the model insufficient information to train on warmer summer days.

The same 20\% of the data was used to evaluate all methods in this study. The OLS regression was fitted on the remaining 80\% and tested on the 20\%. The afternoon classifier, trained on 2,000 synthetic observations, was tested on the same 20\% from the real days. Keeping the test set identical across all methods means the comparisons are fair.

\subsection{Ordinary Least Squares Regression}
\label{sec:ols}

An ordinary least squares (OLS) regression was fitted to the real dataset to quantify the contribution of each demand signal and to provide the statistical foundation for the simulation methodology described in Section~\ref{sec:simulation}. The label is the raw afternoon cover count; the features are hotel occupancy percentage, mean daily temperature, day of week (as an ordinal integer 0--6), the binary special-day indicator, and dummy variables for summer and winter seasons (with autumn as the baseline). Spring was tested and found to be statistically similar to autumn and is discussed below.

The fitted regression equation is:

\begin{equation}
\begin{split}
\widehat{\text{Covers}} = 14.11 &+ 1.73 \cdot \text{occupancy} + 4.46 \cdot \text{temp\_mean} + 7.04 \cdot \text{day\_of\_week} \\
                                &+ 70.67 \cdot \text{is\_special} + 72.81 \cdot \text{Summer} + 85.73 \cdot \text{Winter}
\end{split}
\label{eq:ols}
\end{equation}

\noindent with an overall $F$-statistic of 36.07 ($p < 10^{-32}$) and adjusted $R^2 = 0.418$. Table~\ref{tab:ols} presents the full coefficient table.

\begin{table}[H]
\centering
\small
\begin{tabular}{@{}lrrrr@{}}
\toprule
\textbf{Feature} & \textbf{Coefficient} & \textbf{Std. Error} & \textbf{$t$} & \textbf{$p$-value} \\
\midrule
Intercept             & 14.11  & 23.15  &  0.610  & 0.543 \\
Occupancy (\%)        &  1.73  &  0.305 &  5.660  & $< 0.001$ \\
Mean temperature (°C) &  4.46  &  1.184 &  3.766  & $< 0.001$ \\
Day of week (0-6)    &  7.04  &  2.596 &  2.711  & 0.007 \\
Special day (0/1)     & 70.67  & 23.010 &  3.071  & 0.002 \\
Season: Summer        & 72.81  & 18.837 &  3.866  & $< 0.001$ \\
Season: Winter        & 85.73  & 16.722 &  5.127  & $< 0.001$ \\
Season: Spring        &  9.94  & 17.740 &  0.560  & 0.576 \\
\bottomrule
\multicolumn{5}{@{}l}{\small Adjusted $R^2 = 0.418$; $F(6, 287) = 36.07$, $p < 10^{-32}$.}\\
\multicolumn{5}{@{}l}{\small Spring dropped from final model ($p = 0.576$; indistinguishable from Autumn baseline).}
\end{tabular}
\caption{OLS regression results: afternoon covers on demand signals ($n=294$).}
\label{tab:ols}
\end{table}

Hotel occupancy is the strongest and most precisely estimated predictor. Each percentage point increase in occupancy is associated with approximately 1.73 additional afternoon covers ($t = 5.66$, $p < 0.001$), as reported in \citet{thompson2008}. To put this in context, moving from a 30\% occupancy day to a 70\% occupancy day corresponds to approximately 69 additional predicted covers. This swing is equivalent to roughly a Medium-to-High tier boundary shift.

Temperature, too,  has a significant effect. Each degree Celsius of mean daily temperature adds approximately 4.46 covers ($t = 3.77$, $p < 0.001$). For context, a 15°C temperature difference corresponds to approximately 67 additional predicted covers. This in practice is the difference between a warm summer day at 25°C versus a 10°C day in autumn. Temperature is not merely a stand-in for seasons - season dummies already cover that. It carries additional predictive information within seasons. A warm day in October will likely generate more covers than its colder counterpart, even after taking into account autumn classification.


The day-of-the-week coefficient of 7.04 covers per unit ($t = 2.71$, $p = 0.007$) indicates that each successive day of the week is associated with approximately 7 additional afternoon covers, on average. Moving from Monday (coded 0) to Saturday (coded 5) corresponds to a 35-cover increase. While small, it is still meaningful. In a real-world scenario, for instance, 35 additional covers would require one additional server for the duration of the service.

The special-day coefficient of 70.67 covers ($t = 3.07$, $p = 0.002$) is the largest in the model. Hungarian public holidays and the four additional special dates are not equally busy — the 70.67 figure is an average across all of them. The wide standard error (23.01) shows this clearly: some occasions run well above that, others below.

The winter coefficient of 85.73 (\(t = 5.13\), \(p < 0.001\)) is the most surprising result. Winter months are predicted to generate more covers than autumn, even after accounting for occupancy and temperature. This effect is largely driven by December, when the festive period significantly increases the average number of covers. February shows a slight recovery around Valentine's Day, an effect already captured by the special-day flag. January itself is typically quieter.

Spring days are statistically similar to autumn days ($p = 0.576$). Once occupancy and temperature are accounted for, there is no remaining seasonal effect that is specific to spring alone. This null finding suggests that the seasonal demand effect in spring is fully explained away by the occupancy and temperature signals. In spring, both occupancy and temperature are rising, and these rising signals explain the observed demand recovery without a season-specific effect. Spring was therefore dropped from the final model.

The residual analysis (Figure~\ref{fig:residuals}) shows that the regression residuals are approximately normally distributed and centered on zero, with no obvious evidence of heteroscedasticity. The residual standard deviation is 87.1 covers, compared with a mean of 249 covers. This indicates substantial variation around the model's predictions. This remaining uncertainty justifies the stochastic component of the simulation, as no model can fully predict every group booking, walk-in, or last-minute cancellation.

The OLS is therefore an appropriate estimator for this data. The residuals can be treated as a sample of the unpredictable variation in afternoon covers that features cannot explain. This residual structure is exploited directly during synthetic data generation.

\begin{figure}[H]
\centering
\includegraphics[width=\textwidth]{figures/residuals.png}
\caption{  Left: residual histogram — roughly symmetric, centered near zero. Right: 
  residuals against fitted values, no clear pattern.}
\label{fig:residuals}
\end{figure}

\subsection{Simulation Design}
\label{sec:simulation}

With only 294 real covers data, training a classifier directly on the real data is challenging. A three-class classifier with six features would have approximately 98 observations per class, which may be insufficient for reliable model training \citep{lu2023}. To address this limitation, the simulation generates 2,000 synthetic training data points based on the real data and validates their distribution against the real data before training the classifier.

The simulation was developed through five design stages, each addressing a limitation identified in the preceding stage. This progression is important because a model is only as good as the data it was trained on.

\subsubsection{Stage 1  --  Domain-Knowledge Parameters}

The initial simulation used hardcoded seasonal base rates derived from personal operational experience. For example, summer was assigned a base cover rate of 0.75. These captured intuition based off of experience but had no real data backing. The fundamental problem was that parameters were invented, used in simulation, and partially validated against real data. The right approach is to derive parameters from real data first and use them in simulation.

\subsubsection{Stage 2  --  Statistically Derived Parameters}

All simulation parameters were recomputed from the 294 real cover observations. Monthly base cover rates were estimated as the empirical mean cover rate for each calendar month. Weekend deltas were computed as the mean(weekend) $-$ mean(weekday) difference within each season. The special-day delta was estimated as the mean difference between special and non-special days across all 294 observations ($+0.14$). The noise standard deviation was set to 0.21, computed as the residual standard deviation within season-by-day-of-week cells in the real data.

Table~\ref{tab:monthly} presents the empirically derived monthly base rates.

\begin{table}[H]
\centering
\small
\begin{tabular}{@{}lrrl@{}}
\toprule
\textbf{Month} & \textbf{$n$ days} & \textbf{Base rate} & \textbf{Note} \\
\midrule
January   & 28 & 0.41 & Individual mean \\
February  & 25 & 0.41 & Individual mean \\
March     & 29 & 0.40 & Individual mean \\
April     &  8 & 0.43 & Spring mean substituted ($n$ too small) \\
May       &  7 & 0.43 & Spring mean substituted ($n$ too small) \\
June      & 27 & 0.80 & Individual mean \\
July      & 28 & 0.74 & Individual mean \\
August    & 29 & 0.75 & Individual mean \\
September & 26 & 0.64 & Individual mean \\
October   & 28 & 0.43 & Individual mean \\
November  & 29 & 0.35 & Individual mean \\
December  & 30 & \textit{ramp} & Empirical linear ramp (Stage 4) \\
\bottomrule
\multicolumn{4}{@{}l}{\small Spring mean = 0.43 (computed over all 44 Mar+Apr+May observations).}
\end{tabular}
\caption{Monthly base cover rates derived from 294 real observations.}
\label{tab:monthly}

\end{table}

April and May presented a specific challenge. With only 8 and 7 observations respectively, the 95\% confidence interval on each individual mean spans approximately $\pm 0.20$ -- as wide as the estimate itself. Both months were replaced with the spring season mean (0.43), computed across all 44 observations from March, April, and May combined. This is a known limitation meaning that if April cover rates differ systematically from the spring mean, the simulation will carry a small bias for those months.

\subsubsection{Stage 3  --  Visual Validation: Seasonal Alignment}

  Both series were plotted on the same seasonal window, that is, June to April,  using 
  rolling means, so the shapes could be compared directly.

  The match is strong through most of the year. The summer plateau shows up clearly in both. The autumn drop happens at roughly the same pace, and the January dip lands at about the same level. December had the problem in both graphs. The real data shows a sharp festive recovery that the simulated version, using a flat monthly mean, only reached about 65–70\% of. That gap motivated Stage 4.

\begin{figure}[H]
    \centering
    \includegraphics[width=\textwidth]{figures/seasonal_real_v_simulated.png}
    \caption{Real vs. simulated covers aligned to the same seasonal window (June–April). }
    \label{fig:seasonal-real-simulated}
\end{figure}


\subsubsection{Stage 4  --  December Festive Ramp}

Real covers data reveals that the first two weeks of December are relatively quiet, with a mean cover rate of 0.57, while the final two weeks climb sharply to 0.77 as the festive season takes hold. This is a difference of 0.20 cover-rate units. Applying a single flat mean across the entire month would significantly underestimate the late-December demand. To account for this, December is modeled as a linear ramp between its two empirically observed halves:


\begin{equation}
\text{base}(d) = \mu_{\text{early}} + \frac{d - 1}{30} \cdot (\mu_{\text{late}} - \mu_{\text{early}})
\label{eq:dec_ramp}
\end{equation}

\noindent where $d \in \{1, \ldots, 31\}$ is the day of the month, and $\mu_{\text{early}}$, $\mu_{\text{late}}$ are the empirical means for the early and late halves respectively. 

\subsubsection{Stage 5  --  Multi-Year Cycle Validation}

Figure~\ref{fig:fullrange} confirms that the empirically derived features produce consistent seasonal patterns across multiple years. The simulated graph shows three distinct summer peaks, two winter dips, each reproducing the same shape as the real data's single year. The pattern repeats across three years, which suggests the parameters reflect something real rather than a quirk of the one year available.



\begin{figure}[H]
    \centering
    \includegraphics[width=\textwidth]{figures/full_range_real_vs_sim.png}
    \caption{Real data (Jan 2025–Apr 2026) vs Three years of simulated data. The seasonal structure repeats consistently across  years.}
    \label{fig:fullrange}
\end{figure}

\newpage

Table~\ref{tab:lineage} presents the complete data lineage for all features used in the simulation.

\begin{table}[H]
\centering
\small
\begin{tabular}{@{}p{3.8cm}p{4.5cm}p{1.8cm}@{}}
\toprule
\textbf{Feature} & \textbf{Source} & \textbf{Type} \\
\midrule
\texttt{hotel\_occupancy\_pct}           & 119,390 real hotel bookings        & Real \\
\texttt{is\_special\_day}                 & Hungarian public holiday calendar  & Real \\
\texttt{is\_weekend}, \texttt{day\_of\_week} & Derived from real booking dates & Real \\
\texttt{temp\_mean}, \texttt{temp\_max}   & Open-Meteo Budapest historical API & Real \\
Cover rate: monthly base                 & Mean per month, 294 real obs.      & Empirical \\
Cover rate: weekend delta                & Season-level weekend/weekday diff. & Empirical \\
Cover rate: special day delta (+0.14)    & Mean diff.\ on/off special days    & Empirical \\
Cover rate: December ramp                & Early/late Dec means, real data    & Empirical \\
Cover rate: noise std (0.21)             & Residual std, real data            & Empirical \\
\texttt{has\_local\_event} (+0.10)       & No event log available             & \textbf{Assumption} \\
\bottomrule
\end{tabular}
\caption{Data lineage for simulation features.}
\label{tab:lineage}
\end{table}

The only assumption in the simulation is \texttt{has\_local\_event} (\(+0.10\)). As no event log was available, this effect could not be estimated from the data. The value was therefore set conservatively, at a magnitude below the noise standard deviation of 0.21. This is the only simulation parameter that was not empirically grounded in the real observations.

\subsection{Synthetic Data Generation}
\label{sec:synth}

The synthetic data generation pipeline produces 2,000 training data points. 2,000 data points yield approximately 650--700 observations per class which is sufficient for model training. 2,000 also represents approximately a 6.8$\times$ extrapolation of the real data. Higher amplification risks stretching the regression's domain into regions not well represented in the real data. The synthetic data  points are generated in the following steps:

\paragraph{Step 1: Row sampling:} Feature vectors are generated by sampling entire rows, with replacement, from the 294 real observations. Sampling complete rows preserves the relationships between features as they occurred in the original data. Sampling individual columns independently could produce feature combinations that were never observed in reality.

\paragraph{Step 2: OLS prediction:} Each sampled feature vector is passed through the fitted OLS regression (Equation~\ref{eq:ols}) to obtain a predicted cover count.

\paragraph{Step 3: Residual injection:} A residual is sampled independently from the 294 real OLS residuals and added to the predicted cover count. This independence is intentional. Pairing a feature vector with the residual from its original observation would simply reproduce the observed cover count, rather than generate a new synthetic observation.

\paragraph{Step 4: Clipping and tier labelling:} Cover counts are clipped at zero to avoid negative values. They are then assigned to one of three demand tiers using tertile thresholds from the real data, as in Table~\ref{tab:tier_counts}.

\begin{table}[H]
  \centering
  \begin{tabular}{@{}ll@{}}
  \toprule
  \textbf{Tier} & \textbf{Cover range} \\
  \midrule
  Low    & $\leq 182$   \\
  Medium & $183 - 305$  \\
  High   & $> 305$      \\
  \bottomrule
  \end{tabular}
  \caption{Demand tier thresholds.}
  \label{tab:tier_counts}
  \end{table}

\begin{figure}[H]
\centering
\begin{tikzpicture}[
    node distance=0.55cm,
    box/.style={
        rectangle, rounded corners=4pt, draw=black!60, fill=blue!7,
        text width=10cm, minimum height=1cm, align=center, font=\small
    },
    endbox/.style={
        rectangle, rounded corners=4pt, draw=black!60, fill=gray!12,
        text width=10cm, minimum height=1cm, align=center, font=\small\bfseries
    },
    arr/.style={-{Stealth}, thick}
]
\node[endbox] (src)  {294 Real Afternoon Cover Observations};
\node[box, below=of src] (s1)
    {\textbf{Step 1 — Row Sampling}\\[2pt]
     Sample entire rows with replacement to preserve joint feature correlations};
\node[box, below=of s1] (s2)
    {\textbf{Step 2 — OLS Prediction}\\[2pt]
     Apply fitted regression (Eq.~\ref{eq:ols}) to obtain the predicted cover count};
\node[box, below=of s2] (s3)
    {\textbf{Step 3 — Residual Injection}\\[2pt]
     Add an independently sampled real residual to capture unpredictable variability};
\node[box, below=of s3] (s4)
    {\textbf{Step 4 — Clip and Label}\\[2pt]
     Clip to zero; assign \textit{Low} (${\leq}182$), \textit{Medium} (183--305), or \textit{High} ($>305$)};
\node[endbox, below=of s4] (out)
    {2,000 Synthetic Training Observations (KS-validated, $p = 0.408$)};
\draw[arr] (src) -- (s1);
\draw[arr] (s1) -- (s2);
\draw[arr] (s2) -- (s3);
\draw[arr] (s3) -- (s4);
\draw[arr] (s4) -- (out);
\end{tikzpicture}
\caption{Synthetic data generation pipeline.}
\label{fig:synth_pipeline}
\end{figure}


\subsubsection{Alternative Approaches Considered}


Independent column sampling was considered initially but rejected because it destroys the correlations between features present in the real data. For example, hotel occupancy and summer season are positively correlated (\(r \approx 0.44\)), reflecting the influence of Budapest's tourism season on both. Sampling these features independently would generate unrealistic combinations, such as high occupancy during winter, at frequencies not observed in the original data. Since the OLS regression was fitted on the real joint distribution of the features, its predictions for these artificial combinations could be unreliable. This approach was therefore rejected in favour of sampling complete rows.

The second alternative was \textit{Gaussian parametric simulation}. Here, a multivariate Gaussian distribution is fitted to the feature space and synthetic feature vectors are sampled from it before being passed through the OLS regression. Although this approach preserves linear correlations, it cannot adequately capture non-linear relationships between features. For example, the interaction between occupancy and season is not well represented by a Gaussian distribution, particularly because season is categorical and occupancy is right-skewed. Furthermore, Gaussian sampling can generate values outside the observed range, such as occupancy above 100\%, non-integer day-of-week values, and non-binary special-day indicators. Clipping and rounding these values after sampling would introduce additional distortions. 


Row sampling avoids all this. Every synthetic feature vector is a row in the real data without any extra steps in between.

\subsubsection{Distributional Validation: Kolmogorov--Smirnov Test}

The Kolmogorov--Smirnov (KS) test statistic is defined as:

$$D = \sup_x |F_1(x) - F_2(x)|$$

where $F_1$ and $F_2$ are the empirical Cumulative Distribution Functions of the real and synthetic cover distributions.

\begin{table}[H]
\centering
\begin{tabular}{@{}lll@{}}
  \toprule
   & \textbf{Real} & \textbf{Synthetic} \\
  \midrule
  $n$       & 294   & 2{,}000 \\
  Mean      & 249.4 & 250.7   \\
  $\sigma$  & 116.7 & 116.3   \\
  \midrule
  \multicolumn{3}{@{}l}{KS test: $D = 0.058$, $p = 0.408$} \\
  \bottomrule
  \end{tabular}
  \caption{Real vs.\ synthetic cover distribution comparison.}
  \label{tab:kstest}
  \end{table}

Table~\ref{tab:kstest} shows that the real and synthetic distributions are statistically similar.

\begin{figure} [H]
    \centering
    \includegraphics[width=\textwidth]{sim_vs_real.png}
    \caption{Density histograms of afternoon covers for the real data ($n = 294$) and the synthetic data ($n = 2{,}000$).}
    \label{fig:ksdist}
\end{figure}



\newpage
% =============================================================================
% SECTION 4 — DEMAND TIER CLASSIFICATION MODEL
% =============================================================================
\section{Demand Tier Classification Model}
\label{sec:model}

This section covers the model's development, evaluation, and integration into the scheduling engine.

\subsection{Problem Formulation}


Given a feature vector describing the demand signals on a particular day, the model predicts one of three demand tiers: \textit{Low}, \textit{Medium}, or \textit{High}. Each of these tiers maps to a specific staffing number per role that is encoded in the staffing configuration.
 

Classification works better than regression here for a few reasons. First, tier labels are harder to corrupt through noise. A Saturday that is genuinely busy is unlikely to be labeled 'Low' from random fluctuation alone, as long as the quantile boundaries are set sensibly. And the tier system fits into what was already built, almost naturally. The scheduler already expects a \texttt{role\_count} per role per shift. The \texttt{StaffingService} also maps tiers to those counts through a mutable configuration table. No extra translation layer sits between the model output and the business logic. The three tier output also makes the system easier to use in practice. For instance, if the model predicts high but the manager knows that a large group canceled that morning, the override is very obvious. There is no need to interpret a regression coefficient as would have been in regression.

Tier boundaries were derived from tertile thresholds on the 294 real cover observations. Low is used for days with up to 182 afternoon covers, Medium for 182–305, and High for anything above 305. Tertiles were chosen to keep class sizes roughly even — 97 Low, 101 Medium, 96 High — since imbalanced training data tends to bias classifiers toward the majority class.

Morning staffing is handled differently. Rather than a trained classifier, a direct occupancy threshold rule was applied. That is, below 35\% occupancy maps to Low, 36–65\% to Medium, and above 65\% to High. This came out of the EDA and OLS regression, where hotel occupancy consistently emerged as the variable doing most of the work. Breakfast demand is almost entirely tied to how many guests are in the building that night. A more complex model did not seem warranted when one variable already explained most of what was happening.

\subsection{Feature Set}

Six features are used in the afternoon model, each showing a statistically significant relationship with afternoon cover counts ($p < 0.05$) in the OLS regression:

\begin{enumerate}
    \item \textbf{\texttt{occupancy\_\%}}: Hotel room occupancy percentage (continuous, 0--100).
    \item \textbf{\texttt{temp\_mean}}: Mean daily temperature in degrees Celsius.
    \item \textbf{\texttt{day\_of\_week}}: Day of the week as an integer (0 = Monday, 6 = Sunday).
    \item \textbf{\texttt{is\_special\_day}}: Binary flag (1 if a Hungarian public holiday or operationally significant date).
    \item \textbf{\texttt{season\_Summer}}: Binary dummy (1 if Summer: June--August).
    \item \textbf{\texttt{season\_Winter}}: Binary dummy (1 if Winter: December--February; Autumn is the baseline).
\end{enumerate}

Season Spring was tested in the OLS regression and found to be statistically indistinguishable from Autumn ($p = 0.576$) and was excluded. The feature set thus represents the minimum sufficient set of signals that the data supports.

\subsection{Dataset and Training Split}

The training set contains 2,000 synthetic observations from Section 3.5. The test set is drawn from real observations using a season-stratified split. This means that the last 20\% of each season's days are used, giving 56 real test days in total: Spring (8), Summer (16), Autumn (16), Winter (16).

A naive chronological split would have landed entirely in winter and early spring. That would make it impossible to test the model on High-demand summer days, which is the more interesting evaluation. Splitting by season keeps the temporal order within each season while making sure all seasons are actually represented.

\subsection{Candidate Models}

Four different models are trained on the same 2,000 synthetic training data and evaluated on the same 56 test days. Keeping the data identical across all four validates that performance difference comes from the models themselves, not the data they saw.

\subsubsection{Rule-Based Baseline}

The simplest baseline takes the same occupancy threshold rule used for the morning staffing and applies it to the afternoon staffing. The demand signals, albeit present, are not used. The performance of the system then completely depends on how well occupancy alone can predict afternoon demand. This is the most naive of all the models used. The other models need to beat its accuracy otherwise there is no point in training them.

\subsubsection{Logistic Regression}

A logistic regression classifier is fitted with balanced class weights on standardized features (zero mean, unit variance). The OLS regression already demonstrated that cover volume is approximately linearly related to the six input features. A linear classifier trained on the same feature set is therefore a natural comparator. It tests whether the demand-tier boundary is well-approximated by a hyperplane in the feature space, or whether non-linear methods are necessary.

\subsubsection{Random Forest}

A Random Forest classifier \citep{breiman2001} is trained on the 2,000 synthetic observations, consisting of 200 decorrelated decision trees with balanced class weights. Random Forest does not require feature scaling. It produces feature importance rankings that can be directly compared against the OLS regression coefficients, providing a cross-method validation of which signals drive demand.

\subsubsection{XGBoost}

XGBoost \citep{chen2016} is a gradient-boosted tree ensemble that constructs trees sequentially, each correcting the residual errors of its predecessors. It uses the same 200-estimator budget as the Random Forest to ensure a fair comparison of algorithmic approach rather than compute allocation.

\subsection{Model Comparison and Results}
\label{sec:results}
Table~\ref{tab:modelcomp} summarizes the performance of all four approaches on the 56 real held-out test days.

\begin{table}[H]
\centering
\caption{Comparison of all four models on 56 real test days (season-stratified split).}
\label{tab:modelcomp}
\small
\begin{tabular}{@{}lccccc@{}}
\toprule
\textbf{Model} & \textbf{Accuracy} & \textbf{Low F1} & \textbf{Med F1} & \textbf{High F1} & \textbf{Macro F1} \\
\midrule
Rule-Based           & 44.6\% & 0.35 & 0.50 & 0.44 & 0.43 \\
Logistic Regression  & 62.5\% & 0.59 & 0.38 & 0.89 & 0.62 \\
Random Forest        & 55.4\% & 0.52 & 0.41 & 0.77 & 0.57 \\
XGBoost              & 57.1\% & 0.54 & 0.45 & 0.77 & 0.59 \\
\bottomrule
\multicolumn{6}{@{}l}{\small Season-stratified test split, $n=56$ real afternoon days.}
\end{tabular}
\end{table}

The rule-based baseline achieves 44.6\% accuracy meaning that hotel occupancy alone is a meaningful predictor, but insufficient on its own. The baseline performs about 11\% better than randomness. However, it misclassifies 55.4\% of the days. The full feature set is therefore needed to close this gap.


Logistic Regression achieves the highest overall accuracy at 62.5\% and the highest macro F1 score (0.62), with an F1 of 0.89 on High-demand days. This finding suggests that the demand-tier relationship is largely linear and is consistent since the six features were chosen because of their statistically significant linear relationship with afternoon cover counts. Therefore a linear classifier picking them up well adds up.


Random Forest and XGBoost achieve 55.4\% and 57.1\% accuracy respectively. They both sit squarely above the rule-based baseline but below Logistic Regression. On high demand days, they both score an F1 of 0.77, with a slight difference of 1.7\%. This means that they are picking up on the same patterns. Medium-demand days are the hardest to classify across all four models which tracks with the OLS finding that the boundary between Low and High is the least clear in that range.

The confusion matrices for all four models are shown in Figure~\ref{fig:confmats}:


\begin{figure}[H]
\centering
\includegraphics[width=0.8\textwidth]{two more matrices.png}\\[0.5em]
\includegraphics[width=0.8\textwidth]{two_confusion_matrices.png}
\caption{Confusion matrices for all four models on 56 real test days.}
\label{fig:confmats}
\end{figure}


Logistic Regression is selected as the production model deployed in the integration layer. Its 62.5\% accuracy is the highest across all four approaches. The result is substantively informative: the demand-tier boundary is well-approximated by a linear function of the six input features, consistent with the OLS regression through which the feature set was derived. RF and XGBoost do not improve on this linear baseline given the available training data. This is validation of the feature selection approach.

\newpage
% =============================================================================
% SECTION 5 — SCHEDULER DESIGN AND EVOLUTION
% =============================================================================
\section{Scheduler Design and Evolution}
\label{sec:scheduler}

The SlotMeIn engine went through two versions. A greedy implementation first, then later, a CP-SAT engine. This section covers the formal structure of the problem and how the ML model ties into the final implementation.

\subsection{The Scheduling Problem}

For an assignment to be valid, it must satisfy hard constraints and maximize a weighted objective function. Formally:

$$f: E \times S \rightarrow \{0, 1\}, \quad f(e, s) = x_{e,s}$$

  $$x_{e,s} = \begin{cases} 1 & \text{if employee } e \text{ is assigned to shift } s \\ 0 & \text{otherwise} \end{cases}$$

The task is to find an assignment x that satisfies all hard constraints and maximizes the objective function.

The roles managed by the system are: Team Leader, Server, Bartender, Runner, and Hostess. Each shift slot specifies a required role and a required headcount (\texttt{role\_count}) for that role in that period, determined by the ML-predicted demand tier.

\subsection{Hard Constraints: Labor Law}

The hard constraint set enforces applicable labor regulations unconditionally. No schedule produced by SlotMeIn can violate these constraints, regardless of the demand tier or the configuration of soft preferences.

\paragraph{11-hour rest period:} Article~3 of the EU Working Time Directive \citep{wtd2003} mandates a minimum daily rest period of 11 consecutive hours between working days:

\begin{equation}
\text{start}(s_j) - \text{end}(s_i) \geq 11 \text{ hours} \quad \forall\, e \in E,\ (s_i, s_j) \text{ consecutive assignments of } e
\label{eq:rest}
\end{equation}

\paragraph{Maximum six consecutive working days:} Section~104(3) of the Hungarian Labour Code \citep{hungarian2012} stipulates that employees may not work more than six consecutive days without a rest day. The consecutive-day counter is seeded with the number of consecutive days worked in the preceding week to prevent circumvention at week boundaries:

\begin{equation}
\sum_{d'=d-5}^{d} \mathbf{1}[\text{employee } e \text{ assigned on day } d'] \leq 6 \quad \forall\, e \in E,\ d \in \{1, \ldots, 7\}
\label{eq:consecutive}
\end{equation}

\paragraph{Single shift per day:} Each employee may be assigned to at most one shift slot per calendar day:

\begin{equation}
\sum_{s:\, \text{date}(s) = d} x_{e,s} \leq 1 \quad \forall\, e \in E,\ d \in \{1, \ldots, 7\}
\label{eq:daily}
\end{equation}

\paragraph{Contracted hours:} Each employee has a contracted weekly hours entitlement (full-time: 40 hours; part-time: 30 hours). Total assigned hours must not exceed the contracted entitlement plus a 1-hour operational buffer.

\paragraph{Role eligibility:} An employee may only be assigned to a shift that requires their designated role.

\paragraph{Availability windows:} Assignments are only proposed for shift slots that fall fully within the employee's stated availability window for that day.

\subsection{The Greedy Scheduler: Approach and Limitations}

The first implementation of the SlotMeIn scheduling engine was a greedy constructive heuristic. For each shift slot, a set of eligible employees is generated by applying role eligibility and availability filters. The eligible set is scored using a composite scoring function $S(e) = S_{\text{hours}}(e) + S_{\text{streak}}(e) + S_{\text{rest}}(e)$, where:

\begin{align}
S_{\text{hours}}(e) &= H_{\text{contracted}}(e) - H_{\text{assigned}}(e) \label{eq:score_hours} \\
S_{\text{streak}}(e) &= \begin{cases} 0 & \text{if streak} = 0 \\ -2 & \text{if streak} \in \{1,2\} \\ -5 & \text{if streak} \in \{3,4\} \\ -10 & \text{if streak} \in \{5,6\} \end{cases} \label{eq:score_streak} \\
S_{\text{rest}}(e) &= \begin{cases} -5 & \text{if gap to previous shift} < 11\text{h} \\ 0 & \text{otherwise} \end{cases} \label{eq:score_rest}
\end{align}

Shift slots are processed in ascending order of their eligibility list size. This is the \textit{least-flexible-first} heuristic \citep{rossi2006}, which reduces the risk of reaching unresolvable dead ends.

 The greedy algorithm's fundamental weakness is the permanence of a decision which leaves later, probably harder to fill slots without eligible candidates. This is especially problematic for roles with already few people. For instance, if the restaurant has two Bartenders and the algorithm assigns one on Monday morning,
 suppose the second bartender is unavailable on a more important day like Friday evening, there is no way to undo the Monday assignment and explore alternative options.
 

A related shortfall is ordering sensitivity. Ordering the shift slots by flexibility helps but it does not eliminate the ordering issue. Different orders of the shift slot sequence can produce materially different final schedules from the same employee pool. The order in which the slots are filled determines which employees are exhausted when the next slots are processed. The greedy algorithm is therefore not robust to scheduling horizon changes. Adding a shift slot to the planning week can change the quality of all other slots, not just the new one. This makes the algorithm very brittle. And this brittleness cannot be simply resolved through refining the scoring function.

Lastly, before the integration of the ML component, the staffing levels of the scheduler were determined by a static day-of-the-week lookup. There was no consideration whatsoever for the demand signals. This then means that the schedules produced were identical every week regardless of actual demand. Figures \ref{fig:occupancy} and \ref{fig:greedy_schedules} illustrate this limitation clearly.

\begin{figure}[H]
\centering
\includegraphics[width=0.5\textwidth]{figures/occupancy_18.png}
\caption{The occupancy rates per day, per shift which when used should show a variance in staffing allocation. However the greedy algorithm does not highlight these variances.}
\label{fig:occupancy}
\end{figure}

\begin{figure}[H]
\centering
\includegraphics[width=\textwidth]{figures/greedy_18.png}
\hfill
\includegraphics[width=\textwidth]{figures/greedy_25.png}
\vspace{2em}

\caption{Two consecutive weeks generated by the greedy scheduler; (Oct 18--24 and Oct 25--31). Each role receives the same headcount on the same days in both weeks, with no adjustment for demand signals.}
\label{fig:greedy_schedules}
\end{figure}

Figure~\ref{fig:greedy_schedules} illustrates this limitation concretely. Despite the given occupancy in figure ~\ref{fig:occupancy} being injected to the scheduler, it conveniently fails to use these signals. The schedule then falls short of the actual needs of the given days. Some days like Sunday where occupancy (just one demand signal) is low get proper staffing but this is only by chance.

\subsection{The CP-SAT Scheduler}

The CP-SAT implementation uses Google OR-Tools' CP-SAT solver, effectively replacing greedy with constraint programming. Unlike greedy, it has the ability to backtrack when a partial assignment cannot get the schedule to a valid rendition. It also brings in the ability to optimize multiple objectives at once using a single weighted function.

\subsubsection{Objective Function}

The CP-SAT objective function maximizes a weighted sum over five terms:

\begin{equation}
\max \; \sum_{s \in S} w_{f}\,\mathrm{filled}(s)
  + \sum_{e \in E} \Big[ w_{b}\,\mathrm{constr}(e) - w_{m}\,\delta^{+}_{m}(e) - w_{s}\,\delta^{+}_{s}(e) - w_{u}\,\delta^{-}(e) \Big]
\label{eq:objective}
\end{equation}

where \texttt{filled}(s) is a binary indicator of whether shift slot $s$ is filled; $\delta^+_{\text{mild}}(e)$ and $\delta^+_{\text{steep}}(e)$ are the centihour quantities by which employee $e$'s assigned hours exceed contracted hours in the mild (100--115\%) and steep ($>115\%$) overwork bands; $\delta^-(e)$ is the centihour quantity by which $e$ falls short of contracted hours; and \texttt{constrained}(e) is 1 if employee $e$ is part-time and assigned a shift slot.

Table~\ref{tab:weights} presents the five objective function weights and their operational rationale.

\begin{table}[H]
\centering
\label{tab:weights}
\small
\begin{tabular}{@{}lrrl@{}}
\toprule
\textbf{Term} & \textbf{Weight} & \textbf{Ratio} & \textbf{Rationale} \\
\midrule
\texttt{FILL\_WEIGHT}             & 200 & 1.00 & Unfilled slot is worst outcome; service cannot be delivered \\
\texttt{STEEP\_OVERWORK\_PENALTY} & 100 & 0.50 & Work gets done but at welfare and pay cost \\
\texttt{UNDERWORK\_PENALTY}       &  80 & 0.40 & Financial waste; hours paid for but not utilized \\
\texttt{CONSTRAINED\_FILL\_BONUS} &  50 & 0.25 & Part-time staff are cheaper per hour; margin incentive \\
\texttt{MILD\_OVERWORK\_PENALTY}  &  20 & 0.10 & Expected in hospitality; barely discouraged \\
\bottomrule
\end{tabular}
\caption{CP-SAT objective function weights and operational rationale.}
\label{tab:weights}
\end{table}

The overwork threshold is set at $\Theta = 1.15 \times H_{\text{contracted}}$. Hours up to 115\% of the contracted entitlement incur the mild penalty. Hours above this incur the steep penalty. For a full-time employee working 40 hours/week, the 15\% band represents 6 additional hours before the steep penalty activates. This is well within the EU Working Time Directive \citep{wtd2003}.

\subsubsection{ML Integration: Occupancy-Aware Scheduling}

With the CP-SAT engine in place, ML integration is relatively straight forward at this point. The flow chart below shows this integration clearly:

 \begin{figure}[H]
  \centering
  \begin{tikzpicture}[
      node distance=0.6cm,
      box/.style={rectangle, rounded corners=4pt, draw=black!60, fill=blue!7,
          text width=9.5cm, minimum height=1cm, align=center, font=\small},
      sidebox/.style={rectangle, rounded corners=4pt, draw=black!40, fill=gray!8,
          text width=4cm, minimum height=1cm, align=center, font=\small},
      endbox/.style={rectangle, rounded corners=4pt, draw=black!60, fill=gray!12,
          text width=9.5cm, minimum height=1cm, align=center, font=\small\bfseries},
      arr/.style={-{Stealth}, thick},
      sidarr/.style={-{Stealth}, thick, draw=black!50}
  ]

  \node[box]    (in)   {\textbf{Input:} date, occupancy\_pct, shift (for each date and shift combination)};
  \node[box,    below=of in]  (pred) {\texttt{predict\_tier(date, occupancy\_pct, shift)} \\ Logistic Regression classifier → Tier: \textit{Low}, \textit{Med}, or \textit{High}};
  \node[box,    below=of pred] (staf) {\texttt{apply\_staffing\_to\_periods(tier)} \\ Returns target \texttt{role\_count} for each role in the shift};
  \node[box,    below=of staf] (csp)  {\textbf{CP-SAT Solver} \\ Maximizes coverage objective subject to availability, contract, and fairness constraints};
  \node[endbox, below=of csp]  (out)  {ML-informed Weekly Schedule};

  \node[sidebox, right=1.2cm of csp] (cons) {Talent availability\\Contract constraints\\Shift templates\\Assignment history};

  \draw[arr] (in)   -- (pred);
  \draw[arr] (pred) -- (staf);
  \draw[arr] (staf) -- (csp);
  \draw[arr] (csp)  -- (out);
  \draw[sidarr] (cons) -- (csp);

  \end{tikzpicture}
  \caption{The ML integration pipeline.}
  \label{fig:ml_csat_pipeline}
  \end{figure}

Now, with the ML-powered ability of the scheduler, given the same occupancy as in ~\ref{fig:occupancy}, the schedule produced is as below:


\begin{figure}[H]
\centering
\includegraphics[width=\textwidth]{figures/18_ml.png}
\caption{The week 18-24th October generated using the ML-powered engine}
\label{fig:mlschedule}
\end{figure}



\clearpage
% =============================================================================
% SECTION 6 — RESULTS AND EVALUATION
% =============================================================================
\section{Results and Evaluation}
\label{sec:results}

\subsection{Schedule Comparison: ML-Powered vs Greedy}

The proper evaluation of the SlotMeIn system is a side-by-side comparison of two schedules for the same operational week. There is a configurable table used to define the number of talents to assign given a certain prediction. To illustrate the differences in scheduling outputs, Figure~\ref{fig:staffing_levels} is used in this section. 

 \begin{figure}[H]
     \centering
     \includegraphics[width=0.7\linewidth]{figures/staffing_levels.png}
     \caption{Staffing levels for different roles}
     \label{fig:staffing_levels}
 \end{figure}


The following is an illustrative example using a prospective date beyond the observed data range, to demonstrate how the system behaves on a real input. On Tuesday 20th October 2026, morning occupancy was 67\% and afternoon occupancy was 78\% as in Figure~\ref{fig:staffing_levels}. The temperature read from Open-Meteo~\citep{openmeteo2023} placed Budapest at 9.8\textdegree C. Given these inputs, \texttt{predict\_tier} returned \texttt{high} for the morning shift and \texttt{low} for the afternoon. As seen in Figure~\ref{fig:ml_18}, the Tuesday morning is filled with headcounts from the configuration's busy numbers for all roles and low for the afternoon.


 \begin{figure}[H]
     \centering
     \includegraphics[width=1.05\textwidth]{figures/ml2_18.png}
     \caption{ML schedule for the week of 18th October, 2026}
     \label{fig:ml_18}
 \end{figure}


 \begin{figure}[H]
     \centering
     \includegraphics[width=1.05\textwidth]{figures/greedy_18.png}
     \caption{Greedy schedule for the week of 18th October, 2026}
     \label{fig:greedy_18}
 \end{figure}

  However, the same date, generated using the greedy scheduler does not account for this variance as seen in Figure~\ref{fig:greedy_18}. The juxtaposition in these two figures, Figure~\ref{fig:greedy_18} and Figure~\ref{fig:ml_18}, provides the primary empirical demonstration of the ML integration's effect. 

  In both schedules, certain days and roles have fewer assigned talents than the required headcount. This points to an understaffing problem in the employee pool rather than a limitation of the engine. The engine is only able to assign talents who are available. In case of unavailability of talents to fill a slot, the scheduler skips assignment for that slot and is flagged for manual review by a manager.


\subsection{Model Performance Summary}

Table~\ref{tab:summary} shows the key performance metrics for the demand tier models.

  \begin{table}[H]
  \centering
  \small
  \begin{tabular}{@{}lllll@{}}
  \toprule
  \textbf{Shift} & \textbf{Model} & \textbf{Baseline Acc.} & \textbf{Logistic Acc.} & \textbf{Test $n$} \\
  \midrule
  Morning   & Occupancy threshold rule & --     & --      & -- \\
  Afternoon & Logistic Regression      & 44.6\% & 62.5\% & 56 real days \\
  \bottomrule
  \multicolumn{5}{@{}l}{\small Season-stratified test split}
  \end{tabular}
  \caption{Summary of demand tier model performance.}
  \label{tab:summary}
  \end{table}

The afternoon test accuracy of 62.5\% is the substantive empirical result. The model has learned to predict which of the three demand tiers is needed for restaurant service. Static rules cannot perform this task since it requires reasoning over the demand signals. While this accuracy is good, it is constrained to the demand features fed into it. The literature supports that addition of features that drive demand would in fact drive this figure higher \citep{heliyon2023}.



\subsection{Simulation Validation Summary}

The synthetic training data was validated both visually and statistically. Visually, the time series in both real and simulated data align across the full seasonal cycle. This is evident in Figure~\ref{fig:fullrange}.

Statistically, the KS test yields $p = 0.408$. This confirms that the synthetic cover distribution is statistically similar to the real distribution. All simulation parameters are traceable to real data, apart from the local event flag contribution ($+0.10$). This is explicitly flagged as an assumption in Table~\ref{tab:lineage}.

\subsection{End-to-End Pipeline Summary}

There are five distinct stages in which SlotMeIn connects raw operational data to a published weekly schedule.


\begin{enumerate}
    \item \textbf{Data assembly:} 
    
    The data comes from different sources. Daily occupancy is computed from hotel booking records. Temperature information is collected from OpenMeteo and holidays are explicitly flagged. These data points are then connected by a common calendar date. 

    \item \textbf{Demand signal learning:} An OLS regression is fitted on 294 real afternoon cover observations. This is done to quantify the relationship between contextual features and cover counts. The coefficients and residuals of the fitted regression are exported for use in simulation. This stage is run once and the fitted model is stored.

    \item \textbf{Synthetic data generation and classifier training:} The OLS regression and real residuals are used to generate 2,000 synthetic training observations. A Logistic Regression classifier is trained on the synthetic data and evaluated on real test data. The trained model is serialized to disk. This stage is run once per model generation cycle.

    \item \textbf{Tier prediction:} At schedule generation time, the manager specifies the planning week's start date and the hotel's expected occupancy percentage. The inference wrapper constructs feature vectors for each date-shift combination in the planning week (14 predictions: 7 days $\times$ 2 shifts) and calls the appropriate serialized classifier to obtain the predicted tier for each. Seasonal features (season dummy variables) and calendar features (day of week, weekend, special day) are computed deterministically from the date.

    \item \textbf{CP-SAT scheduling:} The predicted tiers are passed to the \texttt{StaffingService} which maps each tier to a per-role headcount target (\texttt{role\_count}) using the configured \texttt{staffing\_rules} table. Binary decision variables are then constructed for each eligible talent-shift pair. It enforces the hard labor law constraints described in Section~\ref{sec:scheduler}, and maximizes the weighted objective function. The resulting assignment is stored in the database and made available for review and publishing.
\end{enumerate}


\subsection{System Limitations}
\label{sec:systemlimits}

\paragraph{Morning occupancy thresholds unvalidated:} The morning tier thresholds (35\%, 65\%) were set by domain reasoning rather than derived from observed morning operational data. A future version requires morning cover counts or revenue figures to ground these thresholds empirically.

\paragraph{Local event assumption:} The \texttt{has\_local\_event} contribution of $+0.10$ is the only simulation parameter not derived from real data. Its impact on distributional validity is bounded by the KS test result ($p = 0.408$), but it remains an unverifiable assumption.

\paragraph{Manual occupancy input:} The ML tier predictor requires \texttt{occupancy\_pct} as an input at schedule generation time. Until integration with the hotel's property management system is implemented, the manager must enter the expected occupancy figure manually.

\paragraph{Unvalidated objective weight ratios:} The CP-SAT objective function weights are grounded in operational reasoning, but their ratios have not been validated through manager preference elicitation or sensitivity analysis.

\paragraph{Coverage versus equity trade-off:} The CP-SAT objective function rewards filling shift slots and penalises over/underwork, but does not explicitly penalise variance in shifts-per-week across eligible employees. On low-demand weeks where fewer shifts are available than there are staff members, the solver may concentrate assignments on the most available full-time employees while part-time staff receive no shifts. Addressing this would require a soft fairness penalty that penalises variance in weekly shift counts. Alternatively, a two-phase solve could constrain each employee to at most one shift-day in the first phase before filling remaining gaps in the second.
\newpage
% =============================================================================
% SECTION 7 — CONCLUSION
% =============================================================================
\section{Conclusion}
\label{sec:conclusion}


This study has presented SlotMeIn, a scheduling engine that combines machine learning-based predictive analytics and constraint programming to produce schedules that meet the demand of an F\&B restaurant. The underlying argument is that in operations like F\&B where shifts are the main operational structures used, available demand signals can and should be used to make scheduling decisions. The argument has been made at both the empirical and the architectural levels.


Empirically, the study demonstrates that afternoon demand is significantly influenced by various demand signals. These include:
hotel occupancy ($p < 0.001$), mean temperature ($p < 0.001$), day of the week ($p = 0.007$), special-day status ($p = 0.002$), summer season ($p < 0.001$), and winter season ($p < 0.001$). Altogether these signals explain 43\% of the variance in afternoon cover counts. The finding that spring is statistically similar to autumn ($p = 0.576$) is equally informative as it constrains the model and reduces overfitting.

Architecturally, the study presents a pipeline that takes raw operational data and transforms it to a demand-aware schedule. The pipeline consists of three connected stages. First, simulation, which produces empirically grounded training data. Then, a classification stage which applies Logistic Regression trained on synthetic data. And finally, a scheduling stage that uses CP-SAT to enforce hard constraints within a weighted objective function. The hard constraints are derived from the EU Working Time Directive and the Hungarian Labour Code. The scheduler consumes the ML tier prediction as its demand input.

Managers of a F\&B establishment are the obvious direct beneficiaries of SlotMeIn. A system that tells you ahead of time that Tuesday next week is going to be busy or not removes the cognitive burden that currently sits squarely on the memory and experience of one person.  But the staff benefit too.  Overstaffing sends people home earlier than expected, which means fewer hours for them.  While this might not be the worst thing possible, fewer hours also mean less money lands in your account at the end of the month. Understaffing means a service becomes extremely overwhelming. One of the main reasons employees leave hotels is because of the quality of the work environment. Constantly being either overstaffed or understaffed quickly degrades the work environment and causes a high turnover rate. High turnover is also expensive for a hotel. Every replacement costs time and training \citep{tracey2008}. A schedule that distributes hours as fairly as possible and matches staffing to actual demand makes the job more sustainable for the people doing it, which in turn reduces turnover \citep{pham2022, karatepe2020}.

While not primary, the move from greedy to CP-SAT is also a big contribution. It demonstrates the importance of the backtracking ability of the CP-SAT solver in undoing partial assignments which would make the final schedule suboptimal. (Figures~\ref{fig:ml_18} and \ref{fig:greedy_18}) concretely show that ML integration produces schedules that adjust to demand where static rules were not able to make the same adjustment.

A 62.5\% accuracy  for the model does not sound impressive in isolation. However, the baseline, which is the hotel occupancy alone, performs at 44.6\%. That means that adding demand signals improved the accuracy by about 18\%. This means that on roughly 6 out of 10 days, the manager receives the correct tier and schedules the correct number of employees per role even before the shift starts. The remaining days are not random failures. We can attribute them to the scarcity of training and testing data, and genuinely hard to predict days like when a large group cancels last minute, or when a large group comes to the hotel last minute. 

The data scarcity problem is not unique to this restaurant. Most small hospitality businesses will never accumulate enough operational data to train a classifier directly. The methodology used here --- fitting OLS to the available data, deriving simulation parameters from the regression, generating synthetic training data, and validating it with a KS test --- is transferable to any similar context. Arguably, these statistical methods exist to help in similar situations. All a restaurant needs is 200 days of cover data and a property management system to build a pipeline of their own.  The specific parameters may differ but the underlying approach will not.


%\section{Future Work}
\label{sec:future work}
The most impactful extension of this work would be eliminating the need for simulation by gathering more real labeled data from the hotel. This would be about 12--24 additional months of operation. These would cover all the seasons and their patterns across at least two years and allow the model to learn from explicit patterns between the seasons in at least 2 separate years.


Furthermore, an interface that allows employees to make requests for their leave days would also be useful. Currently, a manager would be required to manually input these requests into the system. This work scales linearly with headcount and becomes burdensome quickly. On the same note, integrating a property management system with SlotMeIn would further automate the generation process. This would eliminate the need for a manager to manually input the occupancy rates per shift per day. Consequently, the entire pipeline would then be automatic.

It would also be impactful to do a sensitivity analysis on the objective weights. This would provide the empirical grounding currently missing from the calibration of the objective weights.
This could involve varying each weight by $\pm 20\%$ and comparing the resulting schedules against manager preference judgements. Involving experts in the field for this specific piece of work would substantially strengthen the defensibility of the objective formulation.

Extending the feature set would be the most useful improvement to the afternoon classifier. There are several features not available for this study, yet they influence demand. These include: local events \citep{lasek2016}, online reviews \citep{wu2024}, city-level tourism arrivals \citep{sun2019}, and competitor spatial density \citep{liu2025}. Each of these is supported by peer-reviewed evidence of improving restaurant or hotel demand forecast accuracy. Incorporating them in a future version would be expected to drive afternoon classification accuracy substantially beyond the current 62.5\%.

Finally, the morning tier thresholds are derived from domain reasoning, not observed cover data. They have not been validated against actual morning revenue or cover records. The local event flag is also an assumption with no event log to ground it. While these are not fatal limitations, they are honest ones. A future version of SlotMeIn that validates morning boundaries and incorporates a real event calendar would be much stronger. That is work left to do.
\newpage

% =============================================================================
% ACKNOWLEDGEMENTS
% =============================================================================
\section*{Acknowledgements}
\addcontentsline{toc}{section}{Acknowledgements}

I would like to express my sincere gratitude to my supervisor, Dr.~Bogacsovics Gerg\H{o}. His guidance, constructive feedback and support throughout the development of this thesis have been key to its success. 

I would also like to sincerely thank Bill Omole, my peer mentor. His programming expertise was foundational in the development of the first greedy version of SlotMeIn, the crux of this study. He helped me to hone software engineering principles that show up throughout SlotMeIn.

I am also grateful to the restaurant whose operational data made this empirical work possible. Their willingness to provide access to their data transformed this project from a theoretical exercise into a grounded, evidence-based contribution. 

I also wish to thank the Faculty of Informatics at the University of Debrecen for the academic environment and resources that supported this research. I extend this too to the open-source communities behind Python, scikit-learn, OR-Tools, and the broader scientific computing ecosystem.

Finally, I am deeply grateful to my family and friends for their support. I especially wish to thank Warren Maua, my brother, whose presence has never been absent, and my sister Daisy Awuor, who answered my every call when I needed her. And to Pascal Adina, Kirkstone Osamba, and Derrick Omuga whose doors have always been open to me. You encouraged me to keep going when everything felt impossible, and I could not have finished this without you.



\newpage

\begin{thebibliography}{99}

\bibitem[Lasek et~al.(2016)]{lasek2016}
Lasek, A., Cercone, N., and Saunders, J. (2016).
\newblock Restaurant sales and customer demand forecasting: Literature survey and categorization of methods.
\newblock In \textit{Smart City 360°}, pages 479--491. Springer, Cham.
\newblock DOI: 10.1007/978-3-319-33681-7\_40.

\bibitem[Wu et~al.(2024)]{wu2024}
Wu, D.C., Zhong, S., Song, H., and Wu, J. (2024).
\newblock Do topic and sentiment matter? Predictive power of online reviews for hotel demand forecasting.
\newblock \textit{International Journal of Hospitality Management}, 120:103750.
\newblock DOI: 10.1016/j.ijhm.2024.103750.

\bibitem[Lee(2018)]{lee2018}
Lee, M. (2018).
\newblock Modeling and forecasting hotel room demand based on advance booking information.
\newblock \textit{Tourism Management}, 66:62--71.
\newblock DOI: 10.1016/j.tourman.2017.11.004.

\bibitem[Sun et~al.(2019)]{sun2019}
Sun, S., Wei, Y., Tsui, K.L., and Wang, S. (2019).
\newblock Forecasting tourist arrivals with machine learning and internet search index.
\newblock \textit{Tourism Management}, 70:1--10.
\newblock DOI: 10.1016/j.tourman.2018.07.010.

\bibitem[Liu et~al.(2025)]{liu2025}
Liu, D., Zhao, W., and Luo, L. (2025).
\newblock Restaurant daily demand prediction: Modeling competition and agglomeration effects using spatial-temporal dual graph attention networks.
\newblock \textit{Journal of Hospitality and Tourism Technology}, 17(2):501--.
\newblock DOI: 10.1108/JHTT-01-2024-0028.

\bibitem[Choi et~al.(2009)]{choi2009}
Choi, K., Hwang, J., and Park, M. (2009).
\newblock Scheduling restaurant workers to minimize labor cost and meet service standards.
\newblock \textit{Cornell Hospitality Quarterly}, 50(2):155--167.
\newblock DOI: 10.1177/1938965509333557.

\bibitem[Bergmeir and Ben{\'i}tez(2012)]{bergmeir2012}
Bergmeir, C. and Ben{\'i}tez, J.~M. (2012).
\newblock On the use of cross-validation for time series predictor evaluation.
\newblock \textit{Information Sciences}, 191:192--213.
\newblock DOI: 10.1016/j.ins.2011.12.028.

\bibitem[Kimes(1989)]{kimes1989}
Kimes, S.~E. (1989).
\newblock Yield management: A tool for capacity-considered service firms.
\newblock \textit{Journal of Operations Management}, 8(4):348--363.
\newblock DOI: 10.1016/0272-6963(89)90035-1.

\bibitem[Zippenfenig(2023)]{openmeteo2023}
Zippenfenig, P. (2023).
\newblock Open-Meteo.com Weather API.
\newblock Zenodo.
\newblock DOI: 10.5281/zenodo.7970649.

\bibitem[Perron and Furnon(2024)]{ortools2024}
Perron, L. and Furnon, V. (2024).
\newblock {OR-Tools}.
\newblock Google. \url{https://developers.google.com/optimization}.

\bibitem[Strobl et~al.(2007)]{strobl2007}
Strobl, C., Boulesteix, A.-L., Zeileis, A., and Hothorn, T. (2007).
\newblock Bias in random forest variable importance measures: Illustrations, sources and a solution.
\newblock \textit{BMC Bioinformatics}, 8(1):25.
\newblock DOI: 10.1186/1471-2105-8-25.

\bibitem[Ben Said and Mouhoub(2024)]{bensaid2024}
Ben Said, A. and Mouhoub, M. (2024).
\newblock Machine learning and constraint programming for efficient healthcare scheduling.
\newblock \textit{arXiv preprint arXiv:2409.07547}.

\bibitem[Breiman(2001)]{breiman2001}
Breiman, L. (2001).
\newblock Random forests.
\newblock \textit{Machine Learning}, 45(1):5--32.

\bibitem[Ernst et~al.(2004)]{ernst2004}
Ernst, A.~T., Jiang, H., Krishnamoorthy, M., and Sier, D. (2004).
\newblock Staff scheduling and rostering: A review of applications, methods and models.
\newblock \textit{European Journal of Operational Research}, 153(1):3--27.

\bibitem[European Parliament(2003)]{wtd2003}
European Parliament and Council (2003).
\newblock Directive 2003/88/{EC} of the European Parliament and of the Council concerning certain aspects of the organisation of working time.
\newblock \textit{Official Journal of the European Union}, L~299:9--19.

\bibitem[Hungarian Parliament(2012)]{hungarian2012}
Hungarian Parliament (2012).
\newblock Act {I} of 2012 on the Labour Code (Munka T\"{o}rv\'{e}nyk\"{o}nyve).
\newblock Budapest: Hungarian Parliament.

\bibitem[Lu et~al.(2023)]{lu2023}
Lu, Y., Shen, M., Wang, H., Wang, X., van Rechem, C., and Wei, W. (2023).
\newblock Machine learning for synthetic data generation: A review.
\newblock \textit{arXiv preprint arXiv:2302.04062}.

\bibitem[Dowlut and Gobin-Rahimbux(2023)]{heliyon2023}
Dowlut, N. and Gobin-Rahimbux, B. (2023).
\newblock Forecasting resort hotel tourism demand using deep learning
  techniques: A systematic literature review.
\newblock \textit{Heliyon}, 9(7):e18385.
\newblock DOI: 10.1016/j.heliyon.2023.e18385.

\bibitem[Kim(2025)]{pmc2025}
Kim, S. (2025).
\newblock Development of an {AI}-based restaurant menu demand prediction model
  utilizing sales and meteorological data.
\newblock \textit{Food Science and Biotechnology}, 34(15):3597--3606.
\newblock DOI: 10.1007/s10068-025-01956-2.

\bibitem[Rossi et~al.(2006)]{rossi2006}
Rossi, F., van Beek, P., and Walsh, T. (2006).
\newblock \textit{Handbook of Constraint Programming}.
\newblock Elsevier, Amsterdam.

\bibitem[Thompson(2004)]{thompson2004}
Thompson, G.~M. (2004).
\newblock Workforce scheduling: A guide for the hospitality industry.
\newblock \textit{Cornell Hospitality Report}, 4(6):7--55.
\newblock Center for Hospitality Research, Cornell University.

\bibitem[Thompson and Killam(2008)]{thompson2008}
Thompson, G.~M. and Killam, E.~D. (2008).
\newblock Forecasting covers in hotel food and beverage outlets.
\newblock \textit{Cornell Hospitality Report}, 8(16):4--11.

\bibitem[Van den Bergh et~al.(2013)]{bergh2013}
Van den Bergh, J., Belien, J., De Bruecker, P., Demeulemeester, E., and De Boeck, L. (2013).
\newblock Personnel scheduling: A literature review.
\newblock \textit{European Journal of Operational Research}, 226(3):367--385.

\bibitem[Fang and Shen(2026)]{fang2026}
Fang, Y. and Shen, B. (2026).
\newblock Dynamic staff scheduling optimization algorithm for hotel management.
\newblock \textit{Scientific Reports}.
\newblock DOI: 10.1038/s41598-026-54487-z.

\bibitem[Chen and Guestrin(2016)]{chen2016}
Chen, T. and Guestrin, C. (2016).
\newblock {XGBoost}: A scalable tree boosting system.
\newblock In \textit{Proceedings of the 22nd ACM SIGKDD International Conference on Knowledge Discovery and Data Mining}, pages 785--794.
\newblock DOI: 10.1145/2939672.2939785.

\bibitem[Pham et~al.(2022)]{pham2022}
Pham, N.T., Tučková, Z., and Phan, Q.P.T. (2022).
\newblock Impact of work stress and job burnout on turnover intentions among hotel employees.
\newblock \textit{International Journal of Environmental Research and Public Health}, 19(15):9724.
\newblock DOI: 10.3390/ijerph19159724.

\bibitem[Karatepe et~al.(2020)]{karatepe2020}
Karatepe, O.M., Ozturk, A., and Kim, T.T. (2020).
\newblock Role stress and turnover intention of front-line hotel employees: The roles of burnout and service climate.
\newblock \textit{Frontiers in Psychology}, 10:3217.
\newblock DOI: 10.3389/fpsyg.2019.03217.

\bibitem[Sparks and Browning(2011)]{sparks2011}
Sparks, B.A. and Browning, V. (2011).
\newblock The impact of online reviews on hotel booking intentions and perception of trust.
\newblock \textit{Tourism Management}, 32(6):1310--1323.
\newblock DOI: 10.1016/j.tourman.2010.12.011.

\bibitem[Hyndman and Athanasopoulos(2018)]{hyndman2018}
Hyndman, R.J. and Athanasopoulos, G. (2018).
\newblock \textit{Forecasting: Principles and Practice}, 2nd edition.
\newblock OTexts: Melbourne, Australia.
\newblock Available at: OTexts.com/fpp2.

\bibitem[Tracey and Hinkin(2008)]{tracey2008}
Tracey, J.B. and Hinkin, T.R. (2008).
\newblock Contextual factors and cost profiles associated with employee turnover.
\newblock \textit{Cornell Hospitality Quarterly}, 49(1):12--27.
\newblock DOI: 10.1177/0010880407310191.

\end{thebibliography}

% =============================================================================
% APPENDIX
% =============================================================================
\newpage
\section*{Appendix}
\appendix

All simulation parameters were derived from 294 real afternoon cover observations spanning January 2025 to April 2026. The derivation procedure is fully reproducible by re-running \texttt{notebooks/demand\_model\.ipynb}. The single exception -- the \texttt{has\_local\_event} contribution of $+0.10$ -- is explicitly documented as an assumption, not a data-derived parameter, for the reasons stated in Section~\ref{sec:simulation}.

\newpage
\section*{Declaration on the Use of Artificial Intelligence Tools}
\addcontentsline{toc}{section}{Declaration on the Use of Artificial Intelligence Tools}

I hereby declare that during the preparation of my thesis I have used the following artificial intelligence-based tool(s).

\bigskip
\noindent\textbf{1. Use of AI tools}

\bigskip
\noindent$\boxtimes$ I used artificial intelligence-based tools during the preparation of my thesis as follows:

\bigskip
\noindent\textbf{Name and version of the AI tool:} Claude Code (claude-sonnet-4-6), Anthropic

\bigskip
\noindent\textbf{Exactly how it was used:}
\begin{itemize}
    \item Language correction, stylistic refinement, and text paraphrasing
    \item Explanation and debugging of program code
    \item Writing out reasoning and commentary for notebook sections
    \item Collecting ideas for literature search
\end{itemize}

\bigskip
\noindent\textbf{2. Declaration of independent work}

\bigskip
\noindent I declare that I take full responsibility for the professional content of my thesis. I acknowledge that the use of artificial intelligence-based tools does not replace independent work; therefore, I have verified all statements, results, interpretations, conclusions, and references included in the thesis.

\bigskip
\noindent I further declare that:
\begin{itemize}
    \item the final content of the thesis was compiled by me,
    \item any AI-generated or modified content has been reviewed, verified, and corrected where necessary,
    \item the original results, interpretations, and conclusions presented in the thesis are my responsibility,
    \item the use of AI tools has been truthfully and fully disclosed in this declaration.
\end{itemize}

\vspace{2cm}
\noindent\begin{tabular}{@{}p{6cm}p{6cm}@{}}
Debrecen, 2026 & \dotfill \\
& Maua Elaine Imani \\
\end{tabular}

\end{document}
